\documentclass[11pt]{letter}

\usepackage[T1]{fontenc}
\usepackage[margin=0.8in]{geometry}
\usepackage{microtype}

\signature{Jonathan Pipping-Gam\'on\\Abraham J. Wyner\\Department of Statistics, University of Pennsylvania}
\address{Department of Statistics\\The Wharton School, University of Pennsylvania\\Philadelphia, PA 19104}
\date{August 25, 2026}

\begin{document}

\begin{letter}{Editors\\The Annals of Applied Statistics}

\opening{Dear Editors,}

Please consider our substantially revised manuscript, ``The Blown Lead Paradox: A Pathwise Calibration Benchmark for Win Probability Forecasts,'' for publication in \emph{The Annals of Applied Statistics}.  This submission follows manuscript AOAS2602-023.

The revision places the sports application at the center of the paper.  It begins with a specific NBA collapse, explains why the eventual loser's maximum win probability is a reverse-conditioned path statistic, and develops only the theory needed to benchmark that quantity.  The empirical analysis now includes complete regular-season NFL and NBA inventories for 2018--2024, explicit feed-coverage accounting, and a fixed-clock calibration analysis connected to the sports-forecasting literature, particularly the calibration-surface approach of Yeh, Rice, and Dubin.  The NFL and NBA are presented as parallel applications rather than as a controlled comparison between sports.

We have also strengthened the inferential foundation.  Games sharing teams are handled with a season-stratified dyadic pigeonhole bootstrap, supported by an exchangeable-array empirical-process argument and a simulation study on the observed schedules.  The paper distinguishes this asymptotic justification from finite-sample exactness and reports the expected conservatism of the sparse NFL schedule.  Cross-season franchise linkage, calendar-block shocks, probability rounding, feed corrections, missingness, threshold conventions, and leave-one-season-out analyses are reported transparently.  The complete rebuild leaves the substantive conclusion intact: the NBA feed has an excess of severe eventual-loser peaks, while the NFL analysis finds no detected departure.

The main article now carries one continuous applied argument: the global pathwise departure, its magnitude at the 0.95 benchmark tail, and the first-crossing mechanism.  Proofs, full calibration surfaces, validation, robustness tables, and finite-schedule simulations are in the supplement.  All reported values are generated from one consolidated results file and one TeX macro file.

Thank you for considering the revised manuscript.

\closing{Sincerely,}

\end{letter}
\end{document}
